% appB with its \chapter turned into a section; generated
\section{A short history of computations on the zeros}\label{app:history}

Every verification in Table~\ref{tab:appB} from Backlund's on rests on an evaluator of $\zeta$ on or near the critical line and a certificate that no zero has been missed (Section~\ref{ch:classical}). The window on which Chapters~\ref{chap:measure} to \ref{chap:heat} measure distances, the $79$ zeros below height $200$, is the range that Backlund verified by hand in 1914 \cite{Backlund14}.

\begin{table}[htbp]
\centering\small
\begin{tabular}{@{}llp{0.5\textwidth}@{}}
\toprule
year & authors & result\\
\midrule
1859 & Riemann & \raggedright the first zeros, in unpublished notes later found by Siegel \cite{Siegel32}\tabularnewline
1903 & Gram & \raggedright the first $15$ zeros; Gram points \cite{Gram03}\tabularnewline
1914 & Backlund & \raggedright all $79$ zeros below height $200$ \cite{Backlund14,Backlund}\tabularnewline
1925 & Hutchinson & \raggedright all $138$ zeros below height $300$ \cite{Hutchinson25}\tabularnewline
1936 & Titchmarsh & \raggedright the first $1041$ zeros, all below height $1468$, by the Riemann Siegel formula, computed with the help of Comrie \cite{Titchmarsh36}\tabularnewline
1953 & Turing & \raggedright all $1104$ zeros below height $1540$, on the Manchester computer; Turing's method \cite{Turing53}\tabularnewline
1956 & Lehmer & \raggedright the first $25000$ zeros; the close pair near height $7005$ \cite{Lehmer56,Lehmer56b}\tabularnewline
1969 & Rosser, Yohe, Schoenfeld & \raggedright the first $3.5\times10^6$ zeros; Rosser's rule \cite{RYS69}\tabularnewline
1986 & van de Lune, te Riele, Winter & \raggedright the first $1.5\times10^9$ zeros \cite{LRW86}\tabularnewline
1987 & Odlyzko & \raggedright $10^5$ zeros near the $10^{12}$th zero, and the statistics of their spacings \cite{Odlyzko87}\tabularnewline
1988 & Odlyzko, Sch\"onhage & \raggedright fast evaluation of $\zeta$ at many points \cite{OS88}\tabularnewline
2004 & Gourdon & \raggedright the first $10^{13}$ zeros, not interval certified \cite{Gourdon04}\tabularnewline
2021 & Platt, Trudgian & \raggedright RH verified rigorously below height $3\times10^{12}$ \cite{PT21}\tabularnewline
\bottomrule
\end{tabular}
\caption{Landmarks in the computation of the zeros of $\zeta$. The zeros counted in each row were located on the critical line. For the computations up to 1969 see \cite[Chapters 6 to 8]{Edwards}.}
\label{tab:appB}
\end{table}

Where Turing bounds an average of $S$ to certify a count (Section~\ref{sec:ch5:turing}), the zero density identity of Section~\ref{ch:turing} is exact, an integer at every resolution; at the resolution $\frac12$, the basepoint one of Section~\ref{ch:basepoint}, it is Backlund's method with the vertical side moved from $\re s=2$ to $\re s=1$, and it gives the count $N(T)=M_{1/2}(T)+\tau_{1/2}(T)$ with no hypothesis, which Section~\ref{sec:ch5:count} computes at height $1000$. Lehmer's close pair near height $7005$ \cite{Lehmer56} is computed from the P\'olya kernel in Section~\ref{sec:ch5:lehmer} and followed under the de Bruijn Newman flow in Section~\ref{ch:dbn}. The Odlyzko Sch\"onhage algorithm, with Odlyzko's zeros as the reference measure, is one of the evaluators Section~\ref{sec:ch5:hybrid} proposes for the Thorin comparison at large height.

